\subsection{函数解释为广义函数}

命题 6. --- 设 ν 是 U 上的测度。ν 在 D(U) 上的限制是一个广义函数，它为零当且仅当测度 ν 为零。

设 K 是 U 的紧子集。对每个函数 \(\varphi\in\mathcal{C}_{\mathrm{K}}^{\infty}(\mathrm{U})\) ，我们有 \(|\langle\nu,\varphi\rangle|\leqslant p_{0,\mathrm{K}}(\varphi)|\nu|(\mathrm{K})\) ，因而 \(\nu\) 在 \(\mathcal{D}(\mathrm{U})\) 上的限制是连续的。由于 \(\mathcal{D}(\mathrm{U})\) 在 \(\mathcal{K}(\mathrm{U})\) 中稠密，依据 IV, p. 202 的命题 4(a)，\(\nu\) 在 \(\mathcal{D}(\mathrm{U})\) 上的限制为零当且仅当 \(\nu\) 为零。

我们将把 U 上的复测度空间 \(\mathcal{M}(\mathrm{U};\mathbf{C})\) 与 \(\mathcal{D}'(\mathrm{U})\) 的一个子空间等同起来。我们还将把空间 \(\mathrm{L}_{\mathrm{loc}}^{1}(\mathrm{U})\) 与 \(\mathcal{D}'(\mathrm{U})\) 的一个子空间等同起来，所用的映射是由把 \(f\in \mathcal{L}_{\mathrm{loc}}^{1}(\mathrm{U})\) 映为测度 \(f\cdot \mu\) 的映射过渡到商而诱导的映射（INT, V, p. 44, § 5, no. 2，定义 2）。换言之，对 \(f\in \mathrm{L}_{\mathrm{loc}}^{1}(\mathrm{U})\) 与 \(\varphi \in \mathcal{D}(\mathrm{U})\) ，我们有

\[
\langle f, \varphi \rangle = \int_ {\mathrm{U}} f \varphi d \mu .
\]

特别地，对 \(p \in [1, +\infty]\) ，这使我们可以把空间 \(\mathrm{L}^p(\mathrm{U})\) 与 \(\mathcal{D}'(\mathrm{U})\) 的一个子空间等同起来（参见 INT, V, p. 43, § 5, no. 1）。

命题 7. --- 设 \(p \in [1, +\infty]\) 。\(\mathrm{L}^p(\mathrm{U})\) 到 \(\mathcal{D}'(\mathrm{U})\) 中的单射是连续的。

设 \(f \in \mathrm{L}^{p}(\mathrm{U})\) 。设 K 是 U 的紧子集，用 \(\varphi_{K}\) 表示它的特征函数。对每个支集含于 K 的测试函数 \(\varphi\) ，赫尔德不等式给出

\[
| \langle f, \varphi \rangle | = \left| \int_ {\mathrm{U}} f \varphi d \mu \right| \leqslant \mathrm{N} _ {p} (f) \mathrm{N} _ {q} (\varphi) \leqslant \mathrm{N} _ {p} (f) \mathrm{N} _ {q} (\varphi_ {\mathrm{K}}) p _ {0, \mathrm{K}} (\varphi)
\]

其中 q 是 p 的共轭指数。

特别地，可以把 \(\mathcal{D}(\mathrm{U})\) 与 \(\mathcal{D}'(\mathrm{U})\) 的一个子空间等同起来。

命题 8. --- 空间 \(\mathcal{D}(\mathrm{U})\) 在 \(\mathcal{D}'(\mathrm{U})\) 中稠密。

设 \(\lambda\) 是 \(\mathcal{D}'(\mathrm{U})\) 上在 \(\mathcal{D}(\mathrm{U})\) 上为零的线性形式。由于 \(\mathcal{D}(\mathrm{U})\) 是自反的，存在测试函数 \(\varphi \in \mathcal{D}(\mathrm{U})\) ，使得 \(\lambda(f) = \langle f, \varphi \rangle\) 对所有 \(f \in \mathcal{D}'(\mathrm{U})\) 成立。我们得到

\[
0 = \lambda (\overline {{\varphi}}) = \langle \overline {{\varphi}}, \varphi \rangle = \int_ {\mathrm{U}} | \varphi | ^ {2} d \mu ,
\]

因而 \(\varphi = 0\) 。命题由哈恩--巴拿赫定理得出（EVT, II, p. 49，推论 3 (ii)）。

对 \(h \in \mathcal{C}^{\infty}(\mathrm{U})\) ，连续线性映射 \(\varphi \mapsto h\varphi\) （\(\mathcal{D}(\mathrm{U})\) 到自身的）的转置是 \(\mathcal{D}'(\mathrm{U})\) 到自身的连续线性映射，记作 \(f \mapsto hf\) 。这个定义的合理性在于：若 \(f\) 是与测度 \(\nu\) （\(\mathrm{U}\) 上的）相伴的广义函数，则 \(hf\) 与测度 \(h \cdot \nu\) 相伴。事实上，对每个测试函数 \(\varphi \in \mathcal{D}(\mathrm{U})\) ，可以算出

\[
\langle h f, \varphi \rangle = \langle f, h \varphi \rangle = \int_ {\mathrm{U}} h \varphi d \nu = \int_ {\mathrm{U}} \varphi d (h \cdot \nu).
\]

